% 统计图真值来自 source/普通高考/2014/2014新课标1文(河南,河北,山西).pdf 第 1 页、题 18。
% 原卷图为待作答的频率分布直方图坐标模板，未绘制任何柱体；此处保留空网格。
% x 轴标签依次为 0、75、85、95、105、115、125（0 到 75 为原卷断轴后的等距位置）。
% y 轴刻度为 0.002 至 0.040，轴标题为“频率/组距”；网格覆盖全部刻度位置。
\begin{tikzpicture}
\begin{axis}[exam statistical histogram,
  width=10.0cm,
  height=15.053cm,
  xmin=0, xmax=6.6,
  ymin=0, ymax=0.044,
  axis lines=left,
  axis line style={-stealth},
  xtick={0,1,2,3,4,5,6},
  xticklabels={0,75,85,95,105,115,125},
  ytick={0.002,0.004,0.006,0.008,0.010,0.012,0.014,0.016,0.018,0.020,0.022,0.024,0.026,0.028,0.030,0.032,0.034,0.036,0.038,0.040},
  yticklabels={0.002,0.004,0.006,0.008,0.010,0.012,0.014,0.016,0.018,0.020,0.022,0.024,0.026,0.028,0.030,0.032,0.034,0.036,0.038,0.040},
  ytick scale label code/.code={},
  tick style={draw=none},
  grid=major,
  major grid style={draw=gray!58,line width=0.3pt},
  ylabel={},
  clip=false,
]
    \examstatxlabel{质量指标值};
\examstatylabel{\(\dfrac{\text{频率}}{\text{组距}}\)};
% 原卷在 0 与 75 之间使用断轴锯齿，覆盖连续基线后补回锯齿。
\draw[white,line width=1.3pt] (axis cs:0.18,0) -- (axis cs:0.80,0);
\examstatxbreak[line width=0.45pt]{0.49}{0};
\end{axis}
\end{tikzpicture}
